% Radna skripta za učenje teorije seminara.
%
% Objašnjenja, definicije podloge, primjeri, kontrolna pitanja i bibliografija
% izrađeni su uz AI pomoć kao autorov materijal za učenje (vidi
% docs/AI_UPORABA.md). Nisu tekst za predaju i ne prenose se u seminar.
% Iskaz i skicu dokaza kriterija bijega (odjeljak 5), koji su matematički
% rezultat seminara, autor izrađuje sam.
%
% Prevođenje: latexmk -pdf skripta
% ili: pdflatex skripta && bibtex skripta && pdflatex skripta && pdflatex skripta

\documentclass[11pt,a4paper]{article}

\usepackage[T1]{fontenc}
\usepackage[utf8]{inputenc}
\usepackage[croatian]{babel}
\usepackage{lmodern}
\usepackage[margin=2.5cm]{geometry}
\usepackage{amsmath,amssymb,amsthm}
\usepackage{enumitem}
\usepackage{xcolor}
\usepackage{url}
\usepackage[hidelinks]{hyperref}

\theoremstyle{definition}
\newtheorem{definicija}{Definicija}[section]
\newtheorem{primjer}[definicija]{Primjer}
\newtheorem*{napomena}{Napomena}
\theoremstyle{plain}
\newtheorem{teorem}[definicija]{Teorem}
\newtheorem{propozicija}[definicija]{Propozicija}
\newtheorem{lema}[definicija]{Lema}
\newtheorem{korolar}[definicija]{Korolar}

% Uputa za čitanje: gdje i što pročitati.
\newenvironment{citaj}
  {\par\medskip\noindent\textbf{\color{blue!60!black}Pročitaj.}\ }
  {\par\medskip}

% Kontrolna pitanja za samoprovjeru.
\newenvironment{pitanja}
  {\par\noindent\textbf{\color{orange!70!black}Provjeri razumijevanje.}
   \begin{enumerate}[label=\arabic*., nosep]}
  {\end{enumerate}\medskip}

% Mjesto za autorov tekst.
\newcommand{\upisi}[1]{{\color{gray}\itshape [Upiši: #1]}}

% Mjesto u izvoru na koje se tvrdnja oslanja.
\newcommand{\mjesto}[2]{{\small(\cite{#1}, #2)}}

\newcommand{\C}{\mathbb{C}}
\newcommand{\N}{\mathbb{N}}

\title{Mandelbrotov skup\\\large Radna skripta za učenje teorije}
\author{}
\date{}

\begin{document}
\maketitle

\section*{Kako koristiti skriptu}

Skripta je vodič kroz teoriju potrebnu za seminar, poredan od predznanja
prema metodama vizualizacije. Objašnjenja služe za učenje; svaku tvrdnju
koju ćeš koristiti u seminaru pročitaj u navedenom izvoru, zabilježi točno
mjesto u \texttt{docs/IZVORI.md} i u radu je napiši svojim riječima.

Odjeljak~\ref{sec:kriterij} namjerno je prazan: iskaz i skica dokaza kriterija
bijega matematički su rezultat seminara koji izrađuješ sam. Rješenja primjera
nalaze se u dodatku; pokušaj ih najprije sam.

Odjeljak~\ref{sec:udaljenost} obradi samo ako treća metoda uđe u opseg
(odluka do 5.~11.~2026.).

\tableofcontents

% ---------------------------------------------------------------------------
\section{Predznanje: kompleksni brojevi}
\label{sec:kompleksni}

\begin{citaj}
Bilo koji udžbenik analize ili kompleksne analize koji koristiš na studiju:
geometrijska interpretacija, polarni oblik i nejednakost trokuta.
\end{citaj}

Kompleksan broj $z = x + iy$ poistovjećujemo s točkom $(x, y)$ ravnine. Tako
je skup $\C$ ravnina, a podskupovi od $\C$ su likovi koje možemo nacrtati.

\begin{definicija}[Modul]
Modul broja $z = x + iy$ jest $|z| = \sqrt{x^2 + y^2}$, odnosno udaljenost
točke $z$ od ishodišta. Za $z, w \in \C$ broj $|z - w|$ udaljenost je
točaka $z$ i $w$.
\end{definicija}

Svaki $z \neq 0$ može se zapisati u polarnom obliku $z = r(\cos\varphi +
i\sin\varphi)$, gdje je $r = |z|$, a $\varphi$ argument broja $z$. Pri
množenju se moduli množe, a argumenti zbrajaju. Posebno, pri kvadriranju
vrijedi $|z^2| = |z|^2$, a argument se udvostručuje.

\begin{propozicija}[Nejednakost trokuta]
Za sve $z, w \in \C$ vrijedi
\[
  |z + w| \le |z| + |w|
  \qquad\text{i}\qquad
  |z + w| \ge \bigl|\, |z| - |w| \,\bigr|.
\]
\end{propozicija}

Prva nejednakost daje \emph{gornju}, a druga, obratna nejednakost trokuta,
\emph{donju} ogradu za modul zbroja.

\begin{pitanja}
  \item Kako se pri ponovljenom kvadriranju ponašaju brojevi modula manjeg od
        $1$, jednakog $1$ i većeg od $1$?
  \item Kada u prvoj nejednakosti trokuta vrijedi jednakost? Nacrtaj.
  \item Koja od dviju nejednakosti ti treba ako želiš pokazati da je neki
        modul \emph{velik}?
\end{pitanja}

% ---------------------------------------------------------------------------
\section{Iteracija i orbite}
\label{sec:orbite}

\begin{citaj}
Devaney, uvodna poglavlja o iteraciji funkcija i orbitama
\cite{devaney_first_course}; za intuiciju video \cite{numberphile_mandelbrot}.
\end{citaj}

Cijeli seminar vrti se oko jedne familije polinoma: za svaki parametar
$c \in \C$ promatramo
\[
  f_c \colon \C \to \C, \qquad f_c(z) = z^2 + c.
\]
S $f_c^{\,n}$ označavamo $n$-struku kompoziciju: $f_c^{\,0} = \mathrm{id}$,
$f_c^{\,n+1} = f_c \circ f_c^{\,n}$. To nije $n$-ta potencija.

\begin{definicija}[Orbita]
Orbita točke $z_0 \in \C$ pri preslikavanju $f_c$ jest niz
$z_0, z_1, z_2, \dots$ zadan s $z_{n+1} = f_c(z_n)$, odnosno
$z_n = f_c^{\,n}(z_0)$.
\end{definicija}

\begin{definicija}[Fiksna, periodička i predperiodička točka]
Točka $z$ je \emph{fiksna} ako je $f_c(z) = z$, a \emph{periodička} s
periodom $p \ge 1$ ako je $f_c^{\,p}(z) = z$ i $p$ je najmanji takav broj.
Točka je \emph{predperiodička} ako nije periodička, ali je neki član njezine
orbite periodički.
\end{definicija}

\begin{definicija}[Omeđena orbita i bijeg]
Orbita $(z_n)$ je \emph{omeđena} ako postoji $R > 0$ takav da je
$|z_n| \le R$ za sve $n$. Kažemo da orbita \emph{bježi u beskonačnost} ako
$|z_n| \to \infty$.
\end{definicija}

\begin{napomena}
Omeđenost i bijeg nisu \emph{a priori} jedine mogućnosti: neomeđen niz ne mora
težiti u beskonačnost. Za familiju $f_c$ ta se međumogućnost ne pojavljuje,
što je jedna od posljedica kriterija iz odjeljka~\ref{sec:kriterij}.
\end{napomena}

\begin{primjer}[Orbite točke $0$]
\label{pr:orbite}
Za $c = 0, -1, 1, i, -2$ izračunaj prvih pet članova orbite točke $0$ pri
$f_c$. Je li orbita omeđena? Je li $0$ fiksna, periodička ili predperiodička
točka? Rješenje je u dodatku~\ref{dod:rjesenja}.
\end{primjer}

\begin{pitanja}
  \item Zašto je orbita potpuno određena početnom točkom i parametrom $c$?
  \item Može li orbita biti omeđena, a da nije ni periodička ni
        predperiodička? Gdje si to pročitao?
\end{pitanja}

% ---------------------------------------------------------------------------
\section{Ispunjeni Julijev skup i Julijev skup}
\label{sec:julia}

\begin{citaj}
Douady i Hubbard, pogl.~1 i pogl.~3, prop.~3.1 \cite{douady_hubbard_orsay};
slajdovi \cite{rosehulman_julia}; Devaney, poglavlja o Julijevim skupovima
\cite{devaney_first_course}.
\end{citaj}

Do sada je parametar $c$ bio fiksan, a mijenjali smo početnu točku. Skupovi u
ovom odjeljku žive u \emph{dinamičkoj ravnini} varijable $z$; za svaki $c$
dobivamo drugi skup.

\begin{definicija}[Ispunjeni Julijev skup]
\label{def:Kc}
Ispunjeni Julijev skup polinoma $f_c$ jest skup početnih točaka čija orbita
ne bježi u beskonačnost:
\[
  K_c = \{\, z \in \C : f_c^{\,n}(z) \not\to \infty \,\}.
\]
\end{definicija}
\mjesto{douady_hubbard_orsay}{pogl.~1}

\begin{definicija}[Julijev skup]
Julijev skup jest rub ispunjenog Julijeva skupa, $J_c = \partial K_c$.
\end{definicija}

Skup $K_c$ je kompaktan. Intuitivno, $J_c$ je granica između točaka koje
ostaju i točaka koje bježe: u svakoj okolini točke iz $J_c$ ima i jednih i
drugih, pa se tamo ponašanje orbita mijenja naglo.

\begin{teorem}[Fatou, Julia, 1919.]
\label{tm:dihotomija}
Ako $0 \in K_c$, skup $K_c$ je povezan. Ako $0 \notin K_c$, skup $K_c$ je
homeomorfan Cantorovu skupu.
\end{teorem}
\mjesto{douady_hubbard_orsay}{pogl.~1 i pogl.~3, prop.~3.1}

Točka $0$ je istaknuta zato što je to jedina \emph{kritična točka} polinoma
$f_c$, točka u kojoj je $f_c'(z) = 2z$ jednako nuli. Teorem kaže da ponašanje
te jedne orbite određuje oblik cijelog skupa $K_c$.

\begin{pitanja}
  \item Nacrtaj (grubo) $K_0$. Kakvu orbitu ima točka unutar, a kakvu točka
        izvan jediničnog kruga?
  \item Što je kritična točka i zašto je $f_c$ ima točno jednu?
  \item Nalazi li se dokaz teorema~\ref{tm:dihotomija} u razini seminara?
        Što od njega trebaš znati, a što smiješ samo citirati?
\end{pitanja}

% ---------------------------------------------------------------------------
\section{Mandelbrotov skup}
\label{sec:mandelbrot}

\begin{citaj}
Douady i Hubbard, pogl.~1 \cite{douady_hubbard_orsay}; Branner
\cite{branner1989}; za povijest Mandelbrot \cite{mandelbrot1980}.
\end{citaj}

Sada prelazimo u \emph{ravninu parametra}: svaka točka $c$ predstavlja cijeli
polinom $f_c$. Mandelbrotov skup je ,,katalog'' tih polinoma prema obliku
njihovih Julijevih skupova.

\begin{definicija}[Mandelbrotov skup]
\label{def:M}
Mandelbrotov skup jest skup parametara za koje je ispunjeni Julijev skup
povezan:
\[
  M = \{\, c \in \C : K_c \text{ je povezan} \,\}.
\]
\end{definicija}
\mjesto{douady_hubbard_orsay}{pogl.~1}

\begin{propozicija}[Karakterizacija preko orbite točke $0$]
\label{prop:orbita0}
Vrijedi
\[
  M = \{\, c \in \C : f_c^{\,n}(0) \not\to \infty \,\}.
\]
\end{propozicija}

Ova karakterizacija izravno slijedi iz definicije~\ref{def:Kc} i
teorema~\ref{tm:dihotomija}: orbita točke $0$ ne bježi točno onda kada je
$0 \in K_c$, a to je po teoremu točno onda kada je $K_c$ povezan.
Karakterizacija je ključna za računalo: umjesto ispitivanja povezanosti
beskonačnog skupa dovoljno je pratiti \emph{jednu} orbitu.

\begin{propozicija}[Topološka svojstva]
\label{prop:topologija}
Skup $M$ je kompaktan i povezan.
\end{propozicija}
\mjesto{douady_hubbard_orsay}{pogl.~1}

Povezanost skupa $M$ dokazali su Douady i Hubbard; dokaz je izvan razine
seminara, ali je rezultat važan za metodu iz odjeljka~\ref{sec:mariani}.

\begin{napomena}[Mandelbrotova izvorna parametrizacija]
Mandelbrot je 1980.\ promatrao preslikavanja $z \mapsto \lambda z(1 - z)$
\cite{mandelbrot1980}. Douady i Hubbard navode da je familija
$z \mapsto \lambda z + z^2$ konjugirana familiji $f_c$ za
$c = \lambda/2 - \lambda^2/4$ \mjesto{douady_hubbard_orsay}{pogl.~3}.
Svaki kvadratni polinom konjugiran je točno jednom $f_c$, pa se slike u
različitim parametrizacijama razlikuju samo promjenom koordinata.
\end{napomena}

\begin{pitanja}
  \item Pripada li $c = 1$ skupu $M$? A $c = i$? Iskoristi primjer~\ref{pr:orbite}.
  \item Zašto je propozicija~\ref{prop:orbita0} prirodnija za računalo od
        definicije~\ref{def:M}?
  \item Što znači da su dva preslikavanja konjugirana i zašto konjugirana
        preslikavanja imaju ,,istu'' dinamiku?
  \item Provjeri da familija $z \mapsto \lambda z(1-z)$ daje istu vezu s $c$
        kao familija $z \mapsto \lambda z + z^2$. (Uputa: multiplikator fiksne
        točke $0$.)
\end{pitanja}

% ---------------------------------------------------------------------------
\section{Kriterij bijega (matematički rezultat seminara)}
\label{sec:kriterij}

\begin{citaj}
Slajdovi \cite{rosehulman_julia}; Devaney \cite{devaney_first_course};
Branner \cite{branner1989}. Za predaju je potreban primarni izvor u kojem si
sam pronašao iskaz; zabilježi ga ovdje.
\end{citaj}

Ovaj odjeljak izrađuješ sam: on je matematički rezultat seminara.

\begin{lema}
\upisi{pomoćna tvrdnja o rastu modula}
\end{lema}

\begin{teorem}[Kriterij bijega]
\label{tm:kriterij}
\upisi{iskaz}
\end{teorem}

\begin{proof}[Skica dokaza]
\upisi{skica dokaza svojim riječima}
\end{proof}

\begin{korolar}
\upisi{posljedica za položaj skupa $M$ u ravnini}
\end{korolar}

\begin{pitanja}
  \item Koja nejednakost iz odjeljka~\ref{sec:kompleksni} nosi dokaz?
  \item Razlikuje li se argument za slučaj $|c| > 2$ od slučaja $|c| \le 2$?
  \item Zašto je za točku $c = -2$ bitno je li uvjet strog?
  \item Kako iz kriterija slijedi da napomena iz odjeljka~\ref{sec:orbite}
        vrijedi za orbitu točke $0$?
  \item Što kriterij omogućuje u algoritmu, a što \emph{ne} jamči?
\end{pitanja}

% ---------------------------------------------------------------------------
\section{Od teorije do slike: postupak vremena bijega}
\label{sec:escape}

\begin{citaj}
Heiland-Allen, uvod i odjeljak~1 \cite{heilandallen_escape}; Peitgen i Saupe,
pogl.~4 i dodatak~D \cite{peitgen_saupe1988}.
\end{citaj}

Računalo ne može izračunati beskonačnu orbitu. Postupak vremena bijega zato
radi s dvjema konačnim veličinama: pragom $R$ koji opravdava
teorem~\ref{tm:kriterij} i gornjom granicom $N$ broja iteracija.

\begin{definicija}[Vrijeme bijega]
Za parametar $c$, prag $R$ i granicu $N$ vrijeme bijega jest najmanji
$n \le N$ za koji je $|f_c^{\,n}(0)| > R$. Ako takav $n$ ne postoji, kažemo
da točka nije pobjegla u $N$ iteracija.
\end{definicija}

Slika nastaje tako da se pravokutnik ravnine parametra podijeli na rešetku
točaka, po jedna za svaki piksel, i svakoj se izračuna vrijeme bijega.
Točke koje nisu pobjegle obično se boje crno, a ostale bojom koja ovisi o
vremenu bijega.

\begin{napomena}[Ograničenja]
Ako je $R$ odabran u skladu s teoremom~\ref{tm:kriterij}, zaključak
,,pobjegla je'' pouzdan je dokaz da $c \notin M$. Zaključak ,,nije pobjegla u
$N$ iteracija'' to nije: točka može pobjeći tek nakon $N$ koraka. Crni dio
slike zato je aproksimacija skupa $M$ \emph{izvana}, koja se s rastom $N$
smanjuje. Dodatna su ograničenja konačna rezolucija rešetke i aritmetika s
pomičnim zarezom.
\end{napomena}

\begin{pitanja}
  \item Kako izbor $N$ utječe na sliku, a kako na vrijeme izračuna? Gdje se
        na slici razlika najviše vidi?
  \item Zašto broj računskih operacija ovisi o tome koliko je slike crno?
  \item Smije li se umjesto $|z| > R$ provjeravati $|z|^2 > R^2$? Zašto bi
        to bilo korisno?
\end{pitanja}

% ---------------------------------------------------------------------------
\section{Postupak Mariani--Silver}
\label{sec:mariani}

\begin{citaj}
Mariani, odjeljci o algoritmu i njegovim ograničenjima \cite{mariani2021}.
\end{citaj}

Postupak vremena bijega računa svaki piksel posebno, a najskuplji su upravo
crni pikseli, kojima treba svih $N$ iteracija. Postupak Mariani--Silver
pokušava izbjeći računanje velikih jednoličnih područja.

Ideja: za pravokutnik se izračunaju samo pikseli na njegovu rubu. Ako svi
rubni pikseli imaju istu vrijednost, cijela se unutrašnjost popuni tom
vrijednošću bez računanja. Inače se pravokutnik podijeli na manje i postupak
se rekurzivno ponovi, sve dok pravokutnici ne postanu dovoljno mali da se
izračunaju izravno.

\begin{napomena}[Opravdanje i ograničenje]
Ideja se oslanja na topološka svojstva skupa $M$
(propozicija~\ref{prop:topologija} i srodni rezultati Douadyja i Hubbarda):
skup ne može imati ,,otok'' ili ,,rupu'' koji rub pravokutnika ne bi
osjetio. Postupak, međutim, ne provjerava \emph{neprekidni} rub nego konačno
mnogo rubnih piksela, pa uski dio skupa može proći između dvaju uzoraka.
Rezultat zato ne mora biti identičan postupku vremena bijega.
\end{napomena}

\begin{pitanja}
  \item Na koje se točno svojstvo skupa $M$ Mariani poziva? Je li ono dokazano
        za sam skup $M$, za skupove točaka s jednakim vremenom bijega ili ni
        za jedno? Gdje to piše?
  \item Nacrtaj primjer u kojem uski izdanak skupa prolazi kroz pravokutnik,
        a nijedan rubni piksel ga ne pogađa.
  \item Što ćeš mjeriti da pokažeš prednost nad prvim postupkom: vrijeme,
        broj iteracija ili broj izračunatih piksela? Što je poštenija mjera?
  \item Je li ovo različit algoritam ili inačica postupka vremena bijega?
        Obrazloži.
\end{pitanja}

% ---------------------------------------------------------------------------
\section{Opcionalno: procjena vanjske udaljenosti}
\label{sec:udaljenost}

\begin{citaj}
Vepstas \cite{vepstas_derivatives}; Heiland-Allen, dio o Mandelbrotovu skupu
i odjeljak~5.3 \cite{heilandallen_de}; Peitgen i Saupe, pogl.~4
\cite{peitgen_saupe1988}.
\end{citaj}

Postupak vremena bijega za točku izvan skupa kaže samo koliko je brzo
pobjegla. Procjena vanjske udaljenosti uz orbitu računa i njezinu derivaciju
po parametru $c$ te iz njih procjenjuje koliko je točka udaljena od skupa $M$.
Time se tanki izdanci skupa mogu prikazati i kada su uži od piksela.

Za ovu metodu potrebni su derivacija kompleksne funkcije, rekurzija za
derivaciju orbite te rezultat iz kompleksne analize koji opravdava procjenu.
Formule i pretpostavke pročitaj u izvorima.

\begin{pitanja}
  \item Koji teorem stoji iza procjene i je li ona gornja, donja ili samo
        približna ocjena udaljenosti?
  \item Koja se dodatna veličina računa uz orbitu i koliko to povećava cijenu
        jedne iteracije?
  \item Koju prednost pred prvim dvama postupcima možeš pokazati na slici?
\end{pitanja}

% ---------------------------------------------------------------------------
\appendix
\section{Rješenja primjera}
\label{dod:rjesenja}

Orbite točke $0$ iz primjera~\ref{pr:orbite}:

\begin{center}
\renewcommand{\arraystretch}{1.2}
\begin{tabular}{c|l|l}
  $c$ & $z_0, z_1, z_2, z_3, z_4, z_5$ & ponašanje \\ \hline
  $0$  & $0, 0, 0, 0, 0, 0$ & $0$ je fiksna točka; omeđena \\
  $-1$ & $0, -1, 0, -1, 0, -1$ & $0$ je periodička s periodom $2$; omeđena \\
  $1$  & $0, 1, 2, 5, 26, 677$ & bježi u beskonačnost \\
  $i$  & $0, i, -1+i, -i, -1+i, -i$ & $0$ je predperiodička; omeđena \\
  $-2$ & $0, -2, 2, 2, 2, 2$ & $0$ je predperiodička; omeđena \\
\end{tabular}
\end{center}

Za $c = i$: $(-1+i)^2 + i = -2i + i = -i$ i $(-i)^2 + i = -1 + i$, pa se
orbita od $z_2$ nadalje ponavlja s periodom $2$. Za $c = -2$ orbita stiže u
fiksnu točku $2$; parametar $-2$ leži na rubu skupa $M$, što će biti važno
pri odabiru strogog ili nestrogog uvjeta u kriteriju bijega.

% ---------------------------------------------------------------------------
\bibliographystyle{plain}
\bibliography{literatura}

\end{document}
